\documentclass[a4paper,11pt]{article}
\usepackage[utf8]{inputenc}
\usepackage[T1]{fontenc}
\usepackage[french]{babel}
\usepackage[left=2.5cm,right=2.5cm,top=2cm,bottom=2cm]{geometry}
\usepackage{hyperref}
\usepackage{xcolor}
\usepackage{microtype}
\usepackage{lmodern}

\definecolor{primary}{RGB}{0, 51, 102}
\hypersetup{colorlinks=true, urlcolor=primary}

\begin{document}
\noindent \textbf{Lo\"ic Aron MBASSI EWOLO} \\
ENSEA -- Double Dipl\^ome Ing\'enieur \\
\href{mailto:loicaron.mbassiewolo@ensea.fr}{loicaron.mbassiewolo@ensea.fr} \quad (+33) 7 59 52 33 98 \\
\href{https://github.com/Nameless0l}{github.com/Nameless0l} \quad \href{https://mbassiloic.tech}{mbassiloic.tech}

\vspace{1em}
\noindent Cergy, le \today

\vspace{1em}
\noindent \textbf{Objet : Candidature -- Alternance Ingénieur IA Agentique (BPCE SI)}

\vspace{1.5em}
\noindent Madame, Monsieur,

\vspace{0.8em}
\noindent En tant que stagiaire à l'INRIA Saclay, j'ai architecturé un système multi-agents (Google ADK, RAG, LangChain) orchestrant Gemini LLM avec résolution de conflits entre agents spécialisés, et déployé un pipeline NLP (BERT, scoring) en production. Double diplôme ENSEA/Polytechnique Yaoundé en IA et mathématiques appliquées, avec recherche au MILA Montréal sur la généralisation des LLM. BPCE SI, moteur technologique du groupe BPCE, offre un contexte idéal pour déployer des architectures agentiques intelligentes au cœur de services bancaires numériques innovants.

\vspace{0.8em}
\noindent J'ai conçu à l'INRIA une architecture multi-agents coordinant quatre agents spécialisés (météo, ressources, marchés, risques) via Google ADK : raisonnement collaboratif, planification multi-étapes, résolution de conflits d'objectifs. Chaque agent leverages Gemini LLM pour reasoning ; un orchestrateur central alloue tâches et fusionné résultats. RAG enrichit contexte via APIs externes (OpenWeather, données marchés) ; reduction hallucinations via retrieval augmentation. Déploiement FastAPI production avec gestion erreurs, monitoring, tests robustesse. Résultats : recommandations multi-critères intelligentes. Cette approche -- orchestration agentique, coordinations multi-agents, reasoning distribué -- reproduit les défis que BPCE FS adresse pour services bancaires : agrégation données hétérogènes, décisions intelligentes multi-domaines, assistance personnalisée.

\vspace{0.8em}
\noindent Parallèlement, j'ai intégré Gemini Flash/Pro dans Snappy avec RAG chatbots : retrieval contextuel, fine-tuning léger, gestion conversations multi-turn. Recherche MILA sur généralisation LLM (phénomène grokking) approfondit compréhension théorique. Développement backend bancaire (CSC, Fintech, PCI-DSS) démontre appétence pour régulation et sécurité transactionnelle. Open source Laravel API Generator (génération code via LLM) illustre orchestration LLM pour multi-étapes. Je suis disponible dès septembre 2026.

\vspace{0.8em}
\noindent BPCE SI, par ses enjeux de transformation bancaire digitale et ses besoins d'IA agentique intelligent, offre terreau exceptionnel pour consolider ces compétences. Je serais ravi de discuter comment mon profil peut enrichir vos initiatives d'intelligence bancaire distribuée et d'assistance client agentique.

\vspace{1.5em}
\noindent Veuillez agr\'eer, Madame, Monsieur, l'expression de mes salutations distingu\'ees.
\vspace{2em}
\noindent \textbf{Lo\"ic Aron MBASSI EWOLO}

\end{document}
